\chapter{Compact critical sets and deformation topologies}\label{sec:critical-sets}
The continuous pull-tight argument in \cite{MNIndex}, Section 2.8, separates compactness, descent, and preservation of the homotopy class. Throughout this section $M^{n+1}$ is smooth, closed and connected, $n\geq1$. All cycles have mod-two coefficients and lie in the boundary component of Definition \ref{def:minmax-cycles}. Varifolds need not be rectifiable.

\begin{definition}[A mass cap and its stationary subset]\label{def:pulltight-mass-cap}
For $C>0$ let $K_C$ be the space of all nonnegative $n$-varifolds $V$ on $M$ with $\|V\|(M)\leq C$, equipped with the bounded-Lipschitz metric $\mathbf F_V$ on $G_n(M)$. Write
\[
 S_C=\{V\in K_C:\delta V(X)=0\text{ for every smooth vector field }X\}.
\]
Thus $0\in S_C$. This compact ambient space, rather than a space restricted to integral varifolds, is used in \cite{LiAP}, Section 4.2. The first variation is $\delta V(X)=\int\operatorname{div}_P X\,dV(x,P)$.
\end{definition}

\begin{lemma}[Compactness and joint continuity of first variation]\label{lem:pulltight-compact-cap}
The spaces $K_C$ and $S_C$ of Definition \ref{def:pulltight-mass-cap} are compact and metrizable. If $V_j\to V$ weakly in $K_C$ and $X_j\to X$ in $C^1$, then $\delta V_j(X_j)\to\delta V(X)$.
\end{lemma}
\begin{proof}
The Grassmann bundle of a closed manifold is compact and metrizable. Weak compactness of bounded positive measures on a compact space, with the Riesz representation theorem, makes $K_C$ compact. The bounded-Lipschitz metric induces this weak topology. These are the measure-compactness inputs. For fixed smooth $X$, the function $(x,P)\mapsto\operatorname{div}_P X$ is continuous, so its integral is weakly continuous. Further,
\[
 |\delta V_j(X_j-X)|\leq nC\|\nabla(X_j-X)\|_\infty.
\]
This proves joint continuity, first for smooth fields and then for $C^1$ fields by approximation. Each condition $\delta V(X)=0$ is closed; their intersection $S_C$ is therefore compact.
\end{proof}

\begin{lemma}[A minimizing sequence has a compact nonempty critical set]\label{lem:compact-critical-set}
Let $X$ be a finite simplicial complex, $\Pi$ a class in Definition \ref{def:homotopy-width}, and $\Phi_i\in\Pi$ a minimizing sequence with $L=\mathbf L(\Pi)<\infty$. Its critical set $\mathcal C$ is compact and nonempty, and every $V\in\mathcal C$ has mass $L$.
\end{lemma}
\begin{proof}
The maximal masses $m_i$ converge to $L$, hence are uniformly bounded by some $C>0$. Mass is continuous in the $\mathbf F$ topology, so compactness of $X$ gives a maximizing slice for each $i$. Lemma \ref{lem:pulltight-compact-cap} supplies a convergent subsequence of their varifolds, whose limit belongs to $\mathcal C$. The constant test function one shows that every critical limit has mass $L$.
To prove closedness, suppose $V_j\in\mathcal C$ and $V_j\to V$. From a sequence defining $V_j$ choose a slice of index $i_j>\max\{j,i_{j-1}\}$ whose varifold lies within $1/j$ of $V_j$ and whose mass differs from $L$ by less than $1/j$. These slices converge to $V$ and certify that $V\in\mathcal C$. Thus $\mathcal C$ is closed in compact $K_C$.
\end{proof}

\begin{lemma}[Late near-maximal slices concentrate near the critical set]\label{lem:near-critical-slices}
In the setting of Lemma \ref{lem:compact-critical-set}, let $\mathcal U$ be any open neighborhood of $\mathcal C$ in $K_C$. There are $\eta>0$ and $i_0$ such that
\[
 i\geq i_0,\quad \mathbf M(\Phi_i(x))\geq L-\eta
 \quad\Longrightarrow\quad |\Phi_i(x)|\in\mathcal U.
\]
\end{lemma}
\begin{proof}
Otherwise choose $i_j\to\infty$ and $x_j$ with masses at least $L-1/j$ and varifolds outside $\mathcal U$. Their masses are at most $m_{i_j}\to L$. A subsequence converges in $K_C$ to a member of $\mathcal C$, but also to a point of the closed set $K_C\setminus\mathcal U$. This contradicts $\mathcal C\subset\mathcal U$.
\end{proof}

\begin{lemma}[Joint continuity of diffeomorphic pushforward]\label{lem:pulltight-pushforward-continuity}
Let $\varphi_j,\varphi:M\to M$ be smooth diffeomorphisms with $\varphi_j\to\varphi$ in $C^1$. If $V_j\to V$ weakly with bounded mass, then $(\varphi_j)_\#V_j\to\varphi_\#V$ weakly. If boundary cycles $T_j\to T$ in $\mathbf F$, then $(\varphi_j)_\#T_j\to\varphi_\#T$ in $\mathbf F$. The mod-two pushforward, product and homotopy formulas are external current-theory inputs here; the integer-current calculations alone do not establish this coefficient version. Compare the framework in \cite{MNIndex}, Section 2.
\end{lemma}
\begin{proof}
For a continuous test function $f$ on $G_n(M)$, pushforward integrates
\[
 f(\varphi_j(x),D\varphi_j(P))J_P\varphi_j(x)
\]
against $V_j$. These functions converge uniformly to the corresponding expression for $\varphi$. Uniform mass bounds and weak convergence give the first assertion.
For the flat part, split the difference into $(\varphi_j)_\#(T_j-T)$ and $(\varphi_j)_\#T-\varphi_\#T$. A uniform Lipschitz bound gives a uniform flat-norm operator bound for the first term: push forward a decomposition $T_j-T=A+\partial B$ and bound the masses in dimensions $n$ and $n+1$. For the second, join the sufficiently close maps by the short ambient geodesic homotopy $h_j$. Its time derivative is bounded by a constant times $\sup d(\varphi_j,\varphi)$, and its spatial derivative is uniformly bounded. Since $\partial T=0$, the homotopy formula gives
\[
 \mathcal F((\varphi_j)_\#T-\varphi_\#T)
 \leq \mathbf M((h_j)_\#([0,1]\times T))
 \leq A\sup d(\varphi_j,\varphi)\mathbf M(T)\longrightarrow0.
\]
Here the injectivity radius is positive by compactness. Finally a diffeomorphism is injective, so the rectifiable Jacobian formula gives $|\varphi_\#T|=\varphi_\#|T|$ with no cancellation. Combining the varifold and flat conclusions proves $\mathbf F$ convergence. This identity is not asserted for general noninjective maps.
\end{proof}

\begin{example}[A smooth isotopy need not be continuous in mass]\label{ex:translated-slab-mass-distance}
In the unit flat torus $M=\mathbb R^{n+1}/\mathbb Z^{n+1}$ let $T_t$ be the mod-two boundary of the slab $t<x_{n+1}<t+a$, where $0<a<1/2$. Then $t\mapsto T_t$ is $\mathbf F$-continuous, but for all sufficiently small nonzero $t$,
\[
 \mathbf M(T_t-T_0)=4=2\mathbf M(T_0).
\]
This is the obstruction behind the interpolation warning in \cite{LiAP}, Section 4.2.
\end{example}
\begin{proof}
Each boundary is a pair of parallel unit $n$-tori. Translation is a smooth family of diffeomorphisms, so Lemma \ref{lem:pulltight-pushforward-continuity} gives $\mathbf F$ continuity. For sufficiently small nonzero $t$ the four boundary tori at heights $0,a,t,t+a$ are pairwise disjoint. There is no mod-two cancellation, and each has mass one. Thus mass distance stays four. Using a pair of tori ensures that every cycle lies in the boundary component; an individual coordinate torus would not have that property.
\end{proof}
